\section{Arithmetic of the first interaction}
\label{p1:arith:section}

The first correction to the Poisson law comes from two windows whose integer
arguments follow the same rational progression.  The relevant quantity is the
number of square relations, with their affine signs retained.  Counting only
the pairs that support a relation would lose both the size of the correction
and its dependence on the integer length.

Throughout this section, $I_N=[N,2N)\cap\mathbb Z$, $B=L+1$,
$p=2^{-L}$ and $\lambda=Np$, with $|L-\log_2N|\leq C$.
All constants are uniform in this band.  The square-class vector of a positive
integer $n$ is denoted by $v(n)$; its prime coordinates are $v_\ell(n)\bmod2$.
For $V_x=\{x-1,\ldots,x+L-1\}$, let $E_x$ be the space of even subsets of
$V_x$, and let $S_x$ be its image under the square-class map.  Thus $E_x$ has
dimension $L$.  A selection $A\in E_x$ has affine sign
$(-1)^{\ind_{\{x-1\in A\}}}$ for the run-start event $J_x=1$.

For the arithmetic profiles in this section, write
\[
 \rho(I)=|I|L-\dim\sum_{x\in I}S_x,\qquad
 q(I)=\P(J_x=1\text{ for every }x\in I).
\]
Here $\rho$ includes the local kernels. On the regular set it equals
$\rho_{\mathrm{rel}}$; on arbitrary windows the two quantities must be
distinguished. The exact-marginal probability argument uses
$\rho_{\mathrm{rel}}$.
The affine character on the kernel is either trivial or nontrivial.  Accordingly,
with $\eta(I)\in\{0,1\}$ denoting compatibility, finite Fourier inversion gives
\begin{equation}\label{p1:arith:affine}
 q(I)=p^{|I|}\eta(I)2^{\rho(I)}.
\end{equation}
This formula includes incompatible events.  In particular,
$|\eta(I)2^{\rho(I)}-1|\leq2^{\rho(I)}-1$ whenever the individual systems
have full rank.  For a pair, the same inequality holds even before this
restriction, since a zero-dimensional kernel has no compatibility obstruction.

\subsection{The regular set and the cost of its complement}

For $H\geq2$ put
\[
 K_H(n)=\prod_{\substack{\ell>H\ {\rm prime}\\v_\ell(n)\ {\rm odd}}}\ell.
\]
An integer with $K_H(n)=1$ will be called $H$-defective.  When the diameter of
an interval is less than $H$, every other integer in it has a prime coordinate
which occurs nowhere else in the interval.  Such a coordinate excludes that
integer from every square relation.

\begin{proposition}[Regular set]\label{p1:arith:core}
Let
\[
 G_N=\{x\in I_N:\dim S_x=L\},\qquad Z_N^\circ=\sum_{x\in G_N}J_x.
\]
Then $|I_N\setminus G_N|=N^{o(1)}$, and
\begin{equation}\label{p1:arith:coretransfer}
 d_{\rm TV}(\Law(Z_N),\Law(Z_N^\circ))+
 |\E Z_N-\lambda|+|\E Z_N^\circ-\lambda|
 \ll_C N^{-1+o(1)}.
\end{equation}
More generally, for every deterministic $A\subset I_N$ and every choice of
affine targets on its windows,
\begin{equation}\label{p1:arith:marginalsum}
 \sum_{x\in A}|\P(A_xX=b_x)-p|\leq pN^{o(1)}.
\end{equation}
If the fixed moments of $Z_N$ are uniformly bounded, then, for every fixed
integer $j\geq1$ and every $\varepsilon>0$,
\begin{equation}\label{p1:arith:momenttransfer}
 \big|\E Z_N^j-\E(Z_N^\circ)^j\big|
 \ll_{C,j,\varepsilon}N^{-1+\varepsilon}.
\end{equation}
In particular the second factorial cumulants have the same transfer bound.
\end{proposition}

The proof is given in Section~\ref{p14:comp:regular-transfer}.

\subsection{Sparse hosts and weighted nonaligned pairs}

For separated windows, $|x-y|>L$, a prime larger than $B$ occurs with odd
valuation at most once in either window.  Its equation either pins one
vertex to zero or joins two vertices.  After the pinned components are
removed, each remaining component has one binary degree of freedom.  If
$D(x,y)$ counts the defective vertices and $c(x,y)$ the other unpinned
components, then
\begin{equation}\label{p1:arith:graphbound}
 \rho(x,y)\leq D(x,y)+c(x,y),\qquad D(x,y)=o(B).
\end{equation}
The small prime equations and the two block parities can only reduce this
dimension.

A pair is called aligned if, in one orientation, it satisfies
$bx-ay=h$ for coprime positive $a,b\leq B^3$ and
$|h|\leq(a+b)B$.  Using this slightly enlarged definition includes every
alignment produced by two exact units.  Two different reduced slopes cannot
occur for the same dyadic pair once $N$ is large: distinct slopes are separated
by at least $B^{-6}$, whereas the allowed displacement changes the ratio of
the starts by $O(B^{O(1)}/N)$.

\begin{proposition}[Support and residual mass]\label{p1:arith:hosts}
For every $\varepsilon>0$,
\begin{align}
 \#\{(x,y)\in I_N^2:|x-y|>L,\ \rho(x,y)>0\}
   &\ll_{C,\varepsilon}N^{6/5+\varepsilon},
       \label{p1:arith:hostbound}\\
 \sum_{\substack{x,y\in I_N,\ |x-y|>L\\(x,y)\ {\rm nonaligned}}}
       (2^{\rho(x,y)}-1)
   &\ll_{C,\varepsilon}N^{6/5+\varepsilon}.
       \label{p1:arith:nonaligned}
\end{align}
\end{proposition}

The proof is given in Section~\ref{p14:comp:sparse-hosts}.

\subsection{The two leading rational channels}

An exact unit of the reduced channel $a:b$, $a<b$, is a pair of values
$(at,bt)$ in the two windows.  A relative rational relation selects an even
number of these units.  If the first and last units are $t$ and $t+d$, the
number of selections is
\[
 A_1=1,\qquad A_d=2^{d-2}\quad(d\geq2).
\]
Summing over positions in one window gives $L+1-D$ for a selection of
diameter $D$.  Retaining the predecessor sign changes this to
\begin{equation}\label{p1:arith:diametersign}
 g_L(D)=(L-1-D)\ind_{\{0\leq D\leq L\}}.
\end{equation}
Indeed exactly one of the $L+1-D$ placements has its first selected vertex
at the predecessor; its sign changes from $+1$ to $-1$.  The value
$g_L(L)=-1$ is part of the formula.

Define
\begin{align*}
 U_{a,b}(L)&=\sum_{d=1}^{\lfloor L/b\rfloor}
            A_d(L+1-ad)(L+1-bd),\\
 S_{a,b}(L)&=\sum_{d=1}^{\lfloor L/b\rfloor}
            A_d(L-1-ad)(L-1-bd).
\end{align*}
For $L=bm+r$, $0\leq r<b$ and $m\geq1$, both sums are evaluated by
\begin{equation}\label{p1:arith:channelidentity}
 \begin{split}
 &\sum_{d=1}^{m}A_d(L+s-ad)(L+s-bd)\\
 &\quad=2^{m-1}\big[(b-a)(b+r+s)m+(r+s)^2
                   +(a+b)(r+s)+3ab\big]\\
 &\qquad\quad -(a+b)(L+s)-ab,\qquad s\in\{-1,1\}.
 \end{split}
\end{equation}
When $m=0$ the defining sums are empty and equal to zero.
To verify the identity for $m\geq1$, expand the product and use
\[
 \sum A_d=2^{m-1},\quad
 \sum dA_d=(m-1)2^{m-1}+1,\quad
 \sum d^2A_d=(m^2-2m+3)2^{m-1}-1.
\]
In particular, for $m\geq1$,
\begin{align}
 S_{2,3}(L)&=2^{m-1}\big[(r+2)m+r^2+3r+14\big]-5L-1,
       &&L=3m+r,\quad0\leq r<3,\label{p1:arith:S23}\\
 S_{3,4}(L)&=2^{m-1}\big[(r+3)m+r^2+5r+30\big]-7L-5,
       &&L=4m+r,\quad0\leq r<4.\label{p1:arith:S34}
\end{align}

\begin{theorem}[Complete pair profile]\label{p1:arith:pairs}
The sums below are over ordered separated pairs in $I_N$:
\begin{align}
 \sum_{|x-y|>L}(2^{\rho(x,y)}-1)
 &=\frac N3\{U_{2,3}(L)+U_{3,4}(L)\}
      +O_{C,\varepsilon}(N^{6/5+\varepsilon}),\label{p1:arith:positiveprofile}\\
 \sum_{|x-y|>L}(\eta(x,y)2^{\rho(x,y)}-1)
 &=\frac N3\{S_{2,3}(L)+S_{3,4}(L)\}
      +O_{C,\varepsilon}(N^{6/5+\varepsilon}).\label{p1:arith:signedprofile}
\end{align}
Writing $\mu_N=\E Z_N$ and $\alpha_N=\operatorname{Var}(Z_N)-\mu_N$, one has
\begin{equation}\label{p1:arith:cumulantexpansion}
 \alpha_N=\frac{\lambda^2}{3N}\{S_{2,3}(L)+S_{3,4}(L)\}
            +O_{C,\varepsilon}(N^{-4/5+\varepsilon}).
\end{equation}
Thus, with $r_3=L\bmod3$,
\begin{equation}\label{p1:arith:leadingphase}
 \alpha_N\sim
 \frac{r_3+2}{18\,2^{r_3/3}}\lambda^{5/3}L N^{-2/3}.
\end{equation}
After the entire first-channel term is removed, the next term is
\begin{equation}\label{p1:arith:secondphase}
 \alpha_N-\frac{\lambda^2}{3N}S_{2,3}(L)
 \sim\frac{r_4+3}{24\,2^{r_4/4}}\lambda^{7/4}L N^{-3/4},
 \qquad r_4=L\bmod4.
\end{equation}
\end{theorem}

The exact-unit calculation gives the displayed channel sums. Both
orientations of each of the channels $2:3$ and $3:4$ have a translation
interval of total length $N/3+O(B)$. Private prime coordinates identify
their rational kernels; the exceptional normalizations and all remaining
channels have the stated residual mass. Section~\ref{p14:comp:complete-profile}
gives the complete proof, including the endpoint, contact, and
mean-transfer estimates.

The subtraction in \eqref{p1:arith:secondphase} must retain the whole finite
expression $S_{2,3}(L)$.  Its own term of size $N^{-2/3}$ without the factor
$L$ is larger than $LN^{-3/4}$.  The two exact channel polynomials together
depend on $L\bmod12$.

\subsection{Good edges, compression, and one additional window}

Let $\mathcal E_N$ consist of the unordered compatible dependent pairs in
the two channels above for which the distinct normalized values have private
prime coordinates above the fixed threshold $H$.  Write
$q_e=\P(J_x=J_y=1)$ for $e=\{x,y\}$ and
$\beta_N=\sum_{e\in\mathcal E_N}q_e$.  The family is deterministic.

\begin{corollary}[The pair input for the expansion]\label{p1:arith:edgefamily}
One has $|\mathcal E_N|=O_C(NB)$ and, for an edge in channel $a:b$,
\[
 \rho(e)\leq L/b+O(1),\qquad q_e\leq p^2 2^{L/b+O(1)}.
\]
Moreover,
\begin{align}
 2\beta_N&=\alpha_N+O_{C,\varepsilon}(N^{-4/5+\varepsilon}),
    \label{p1:arith:edgemass}\\
 \sum_{\substack{x,y\in G_N,\ x\ne y\\\{x,y\}\notin\mathcal E_N}}
   |\E[(J_x-p)(J_y-p)]|
 &\ll_{C,\varepsilon}N^{-4/5+\varepsilon}+B/N.
    \label{p1:arith:badpairs}
\end{align}
In particular $\beta_N\asymp_C LN^{-2/3}$.
\end{corollary}

The proof is given in Section~\ref{p14:comp:good-edge-mass}.

\begin{lemma}[Exact compression of a good edge]\label{p1:arith:compression}
For a good edge in channel $a:b$, the joint homogeneous event space is the
injective square-class image of the even-selection space of
\[
 U_e=bV_x\cup aV_y,
\]
where repeated values are included only once.  In particular its rank is
\begin{equation}\label{p1:arith:compressedrank}
 \dim(S_x+S_y)=|U_e|-1\geq(2-1/b)L-O(1).
\end{equation}
After the endpoint offsets are fixed, $U_e$ is a subset of an interval of
length $O(B)$ depending on one integer translation parameter.  There are
only $B^{O(1)}$ offset patterns.
\end{lemma}

\begin{proof}
Multiplying every integer in a window by the same constant does not change
the square class of an even selection.  The two even-selection spaces on
$bV_x$ and $aV_y$ therefore map to the original event spaces.  They intersect
on at least one normalized value: dependence in the good family requires
at least two exact units.  The sum of the even-selection spaces of two
intersecting finite sets is the even-selection space of their union.  Its
square-class map is injective, because every distinct value in $U_e$ has a
private prime.  There are at most $L/b+O(1)$ repeated values, which proves
\eqref{p1:arith:compressedrank}.
\end{proof}

For the positional comparison, one also needs the positive conditional
cost of an additional window after a good edge. Proposition~\ref{p1:arith:palm}
in the companion gives this cost separately for the channels $2:3$ and $3:4$.
The optimal scalar reconstruction will instead retain every pair covariance.
